\section*{Chapter \ref{ch:in:how-pay-works} --- How Pay Works}

\begin{solution}{exo:in:how-pay-works:1}
From \$257\,500 to \$180\,000, $-30.1\%$; from \$176\,500 to \$244\,700, $+38.6\%$ on the first estimates. The 2025
release reports $+31.5\%$, because it compares with a revised 2023 figure: estimates are revised, so compare a release
with its own previous-year figure.
\end{solution}

\begin{solution}{exo:in:how-pay-works:2}
40\% of the first \pounds660\,000 is \pounds264\,000 and 60\% of the remaining \pounds340\,000 is \pounds204\,000:
\pounds468\,000, 46.8\% of the bonus.
\end{solution}

\begin{solution}{exo:in:how-pay-works:3}
$1.05^{-4}=0.82$ per unit of award.
\end{solution}

\begin{solution}{exo:in:how-pay-works:4}
The year-1 award vests at the ends of years 2 to 5, so three tranches remain; the year-2 award's four tranches all
remain. Each tranche is $0.4\times200\,000\,e^{0.125}/4=\$22\,663$; seven tranches are \$158\,641, which the simulation
reports as about \$158\,000.
\end{solution}

\begin{solution}{exo:in:how-pay-works:5}
It pays $\max(\text{share}\times\text{P\&L},0)$: a call on the book's result struck at zero. Its payoff is convex, with a
mass at zero and a long right tail; a concave utility weights the zeros heavily, so the certainty equivalent falls
further below the mean as risk aversion rises, as the platform curve shows.
\end{solution}

\begin{solution}{exo:in:how-pay-works:6}
She keeps it with probability 0.85: $0.85\times100\,000=\$85\,000$ in expectation, before discounting and before the
cost of feeling bound to stay.
\end{solution}

\begin{solution}{exo:in:how-pay-works:7}
With the same simulated careers, the expected \omterm{def:in:how-pay-works:rsu}{forfeiture} falls from \$55\,022 to zero and the mean present value rises by
the same amount, from \$1\,617\,971 to \$1\,672\,994: a \omterm{def:in:how-pay-works:rsu}{good-leaver provision} is worth exactly what the leavers would
otherwise have forfeited.
\end{solution}

\begin{solution}{exo:in:how-pay-works:8}
The average covers every securities employee in the city, from operations to senior traders, so it describes no role;
it is a mean over a skewed distribution; and the next year's figure moved by about 30\% in either direction in 2022 and 2024. Chapter~14 gives role and
seniority ranges.
\end{solution}

\begin{solution}{pb:in:how-pay-works:1}
\begin{enumerate}
\item Fixed contractual pay; base plus the year's variable award at award value. It omits \omterm{def:in:how-pay-works:rsu}{forfeiture} risk, the
instrument's risk, timing and taxes.
\item \$257\,500, \$180\,000, \$176\,500, \$244\,700, \$246\,900; about 42\% of wages in 2024.
\item The firm has a fixed cost base and a result that swings with markets; variable pay shares the swing.
\item Formulaic is convex with a mass at zero, bears the book's risk and ends with the book; discretionary is smoother,
bears the pool's risk and the manager's judgement, and can go to zero without the job ending.
\item The sure amount with the same expected utility; it depends on risk aversion and other wealth.
\item A promise of a share on vesting; loss of unvested awards on leaving; a term letting some leavers keep them.
\item EU: at least 40\% deferred over four to five years, pro rata, 60\% for large amounts, half in instruments,
variable at most 100\% (200\%) of fixed. UK: four-year minimum deferral, 40\% on the first \pounds660\,000 and 60\% above.
One US bank: RSUs delivered over three years.
\item Time, the instrument's risk, the probability of leaving (and good-leaver terms), malus and clawback.
\item When it vests, on its market value then.
\item About \$90\,000, \$158\,000, \$203\,000, \$226\,000, \$226\,000; leaving after year two forfeits a present value of
about \$129\,000.
\item A joining payment, repayable if leaving early, which pays for what the candidate gives up and retains her; a fixed
minimum award, which removes the first year's risk.
\item As in the tutorial: bank base 250\,000, median award 200\,000, 40\% deferred into stock over four years; market maker
base 200\,000, median award 250\,000 in cash; platform base 175\,000, 2\% of the book's profit, closed after a loss
beyond 10 million, sign-on 100\,000 and guarantee 150\,000. Five years, 10\% leaving hazard, 5\% discount rate, 500\,000
of wealth.
\item Bank \$1\,618 thousand (597 to 2\,230); market maker \$1\,949 thousand (686 to 2\,956); platform \$1\,703 thousand
(171 to 3\,010).
\item 35.4\%.
\item A share of a fund's profits paid to its managers through the general partner; a member's profit entitlement in a
partnership. UK: income tax on 72.5\% of qualifying \omterm{def:in:how-pay-works:carry}{carried interest} from 6 April 2026; US: a three-year holding period.
\item At relative risk aversion 1: bank \$1\,500 thousand, market maker \$1\,779 thousand, platform \$1\,439 thousand; at
3: \$1\,186, \$1\,365 and \$904 thousand. The buyout for a move from the bank at the end of year two is about \$158\,000
of new awards on the old vesting dates.
\item Its outcomes are the most spread, with a heavy lower tail from closed books; concave utility penalises that most.
\item Its vesting dates, whether it is itself deferred and forfeitable, repayment terms, and good-leaver treatment.
\item The ranking by level depends on the illustrative medians; the shapes (the platform's tail, the bank's handcuff) come
from the structures.
\item Value each offer with her own risk aversion and the real terms, not the headline; if she cannot bear a year of no
pay, the platform is worth well below its mean, and the bank's deferral is a cost to price if she may move.
\end{enumerate}
\end{solution}

\begin{solution}{iq:in:how-pay-works:1}
A call option on the book's annual P\&L struck at zero, with notional 2\%; it is also a short put on her job if a loss
closes the book.
\iqlookfor{convexity and the hidden short option of the loss limit.}
\end{solution}

\begin{solution}{iq:in:how-pay-works:2}
To align pay with the risk taken: deferred pay can be reduced if losses appear later, and shares tie the payee to the
firm's long-run value rather than a single year's result.
\iqlookfor{time-alignment of risk and reward; malus.}
\end{solution}

\begin{solution}{iq:in:how-pay-works:3}
$\sqrt{100\times300}-100=73.2$, against a mean of 100.
\iqlookfor{the geometric mean of wealth, minus wealth.}
\end{solution}

\begin{solution}{iq:in:how-pay-works:4}
One row per tranche: award id, grant date, vest date, units or amount, instrument, conditions (service, performance,
good-leaver classes), status. The unvested balance on a date sums tranches with vest date after it; \omterm{def:in:how-pay-works:rsu}{forfeiture} on a
leaving date applies the conditions to those tranches.
\iqlookfor{tranche granularity and conditions as data.}
\end{solution}

\begin{solution}{iq:in:how-pay-works:5}
Malus reduces an award before it vests; clawback recovers pay already delivered. Clawback is harder, because it must
recover money from a person, possibly a former employee, rather than cancel a book entry.
\iqlookfor{before against after vesting, and enforceability.}
\end{solution}

\begin{solution}{iq:in:how-pay-works:6}
Compare leaving rates of otherwise similar employees with and without unvested balances, or around vesting dates (a
jump in leaving just after vesting is the signature). Data: individual award schedules and leaving dates, with controls
for performance and market conditions.
\iqlookfor{an identification strategy, not an assertion.}
\end{solution}
